\section{Introducción}

\subsection{Motivación y Planteamiento del Problema}

Las arquitecturas de transformers han transformado fundamentalmente el panorama del modelado de secuencias en el procesamiento del lenguaje natural \citep{vaswani_attention_2017}, la visión por computadora y la biología computacional. El éxito de modelos como BERT \citep{devlin_bert_2019}, GPT y sus variantes ha establecido al transformer como el estándar de facto para el aprendizaje de representaciones en el aprendizaje profundo. Sin embargo, la rápida proliferación de innovaciones arquitectónicas presenta desafíos prácticos significativos para investigadores y profesionales.

Los últimos años han sido testigos de una explosión de mecanismos alternativos de atención y mezcla de secuencias, cada uno abordando limitaciones específicas de la atención softmax estándar: complejidad computacional cuadrática \citep{child_sparse_transformer_2019}, requisitos de inferencia eficiente en memoria \citep{sun_retentive_2023,gu_mamba_2023}, gestión selectiva de estados basada en contenido \citep{behrouz_titans_2025} y optimizaciones conscientes del hardware \citep{ramapuram_theory_2024}. Simultáneamente, la literatura sobre optimización se ha diversificado más allá del AdamW clásico, introduciendo técnicas de reducción de varianza \citep{xie_adan_2022,pagliardini_ademamix_2024,yuan_mars_2024}, variantes eficientes en memoria \citep{shazeer_adafactor_2018,zhao_galore_2024}, enfoques sin programación de tasa de aprendizaje \citep{defazio_road_2024} y métodos de precondicionamiento de segundo orden \citep{gupta_shampoo_2018,vyas_soap_2024}.

La consecuencia práctica es un ecosistema de investigación fragmentado donde la comparación experimental entre arquitecturas y optimizadores requiere un esfuerzo de ingeniería significativo. Los investigadores deben implementar y depurar múltiples variantes desde cero, garantizar tuberías de entrenamiento consistentes y gestionar espacios de hiperparámetros complejos. Esta fragmentación dificulta la reproducibilidad, ralentiza el progreso científico y aumenta la barrera de entrada para nuevos investigadores.

Este trabajo aborda estos desafíos mediante un conjunto de herramientas de experimentación unificado y basado en configuración, disponible en \url{https://github.com/erickfmm/frankestein-transformer}, que proporciona:
\begin{enumerate}[leftmargin=*]
    \item \textbf{Diseño Basado en Esquemas}: Un contrato de configuración estricto y validado que garantiza la reproducibilidad mientras admite diecisiete arquitecturas de mezcladores de secuencia distintas y veintitrés familias de optimizadores.
    \item \textbf{Entrenamiento Agnóstico a la Arquitectura}: Infraestructura de entrenamiento común que soporta líneas base de atención densa (estándar, sigmoide), alternativas recurrentes (RetNet, Mamba, bloques tipo EDO), enrutamiento adaptativo de profundidad (Mixture-of-Depths) \citep{zhu2026mixtureofdepthsattention}, bloques de memoria condicional (Engram) \citep{cheng2026conditionalmemoryscalablelookup}, patrones de atención dispersa (Sparse Transformer, Longformer, BigBird, SparseK, NSA, SpargeAttn, FASA) y mecanismos con compuerta (GLA, DeltaNet, Gated DeltaNet, Gated DeltaNet-2, HGRN2, FoX, Gated Softmax).
    \item \textbf{Marco de Enrutamiento de Optimizadores}: Grupos de hiperparámetros con prefijos que permiten un control detallado sobre incrustaciones, capas de normalización, bloques recurrentes, bloques de atención y otros subconjuntos de parámetros a través de diversos optimizadores, incluyendo la familia APOLLO (Apollo, Apollo-Mini, Q-Apollo) \citep{zhu2025apollosgdlikememoryadamwlevel}.
    \item \textbf{Flujos de Trabajo de Extremo a Extremo}: Despliegue integrado mediante cuantización (empaquetado ternario de pesos, activaciones INT8) y capacidades de incrustación de oraciones inspiradas en SBERT \citep{reimers_sentence-bert_2019}.
    \item \textbf{Configuración Interactiva}: Interfaz basada en web que proporciona renderizado de formularios impulsado por esquemas, validación en tiempo real y generación de comandos CLI.
    \item \textbf{Herramientas de Confiabilidad}: Pruebas automatizadas integrales con suites de pruebas unitarias, validación de presets YAML y automatización de CI en versiones de Python compatibles.
\end{enumerate}

La superficie de comandos del proyecto es:
\begin{center}
\texttt{frankestein-transformer}
\end{center}
con subcomandos para flujos de trabajo de codificador y decodificador:
\texttt{train}, \texttt{finetune}, \texttt{deploy}, \texttt{quantize}, \texttt{infer}, \texttt{sbert-train}, \texttt{sbert-infer}, \texttt{web-server}.

El comando \texttt{web-server} lanza un constructor de configuración basado en Streamlit que proporciona:
\begin{itemize}[leftmargin=*]
    \item Campos de formulario impulsados por esquemas con títulos de parámetros y descripciones detalladas
    \item Información sobre herramientas en tiempo real y texto de ayuda para cada opción de configuración
    \item Vista previa YAML en vivo y funcionalidad de descarga
    \item Comandos CLI generados para entrenamiento, despliegue, inferencia y flujos de trabajo SBERT
\end{itemize}

Esta interfaz interactiva sirve como alternativa a la edición manual de YAML, mejorando la usabilidad para usuarios que exploran las opciones de configuración disponibles y comprenden su impacto en el comportamiento del modelo y la dinámica de entrenamiento.

\subsection{Contribuciones}

Este trabajo realiza las siguientes contribuciones principales:
\begin{enumerate}[leftmargin=*]
    \item \textbf{Esquema de Configuración Unificado}: Un contrato de esquema basado en YAML con validación estricta que soporta diecisiete arquitecturas de mezcladores de secuencia distintas en cuatro categorías (líneas base densas, bloques recurrentes/retentivos, patrones de atención dispersa y mecanismos con compuerta), junto con controles de profundidad adaptativa y memoria condicional, mientras garantiza la reproducibilidad mediante restricciones \texttt{additionalProperties: false}.
    \item \textbf{Soporte Integral de Arquitecturas}: Implementación de variantes modernas de transformers incluyendo atención softmax estándar \citep{vaswani_attention_2017}, atención sigmoide \citep{ramapuram_theory_2024}, RetNet \citep{sun_retentive_2023}, Mamba \citep{gu_mamba_2023}, transformers continuos tipo EDO \citep{zhang_continuous_2021}, atención con aumento de memoria Titans \citep{behrouz_titans_2025}, capas de memoria condicional Engram \citep{cheng2026conditionalmemoryscalablelookup}, enrutamiento de tokens Mixture-of-Depths \citep{zhu2026mixtureofdepthsattention}, mecanismos de atención dispersa \citep{child_sparse_transformer_2019,beltagy_longformer_2020,zaheer_bigbird_2020,lou_sparsek_2024,yuan_nsa_2025,zhang_spargeattn_2025,wang_fasa_2026} y arquitecturas de atención con compuerta \citep{yang_gla_2023,yang_deltanet_2024,yang_gated_deltanet_2024,hatamizadeh_gated_deltanet2_2026,qin_hgrn2_2024,lin_forgetting_transformer_2025,qiu_gated_attention_2025}. El sistema soporta tanto entrenamiento en modo codificador con atención bidireccional para modelado de lenguaje enmascarado (MLM) como entrenamiento en modo decodificador con enmascaramiento causal para predicción autorregresiva del siguiente token.
    \item \textbf{Marco de Enrutamiento de Optimizadores}: Sistema de hiperparámetros con prefijos que permite control por grupo de parámetros a través de veintitrés optimizadores, incluyendo Anon con adaptabilidad ajustable \citep{anon2026anon}, métodos de reducción de varianza (MARS, Adan, AdEMAMix, Cautious AdamW), variantes eficientes en memoria (Adafactor, GaLore, Lion), enfoques sin programación de tasa de aprendizaje (Schedule-Free AdamW), métodos conscientes de curvatura (Sophia), precondicionadores de segundo orden (Shampoo, SOAP), optimizadores orientados a ortogonalidad (Muon, Turbo-Muon), escalado de lotes grandes (LAMB) y la familia APOLLO (Apollo, Apollo-Mini, Q-Apollo) \citep{zhu2025apollosgdlikememoryadamwlevel}.
    \item \textbf{Cuantización y Despliegue}: Tubería de despliegue integrada que soporta empaquetado ternario de pesos y cuantización de activaciones INT8 con estimaciones de tamaño siguiendo la aproximación de almacenamiento de $1.58$ bits.
    \item \textbf{Flujos de Trabajo de Incrustación de Oraciones}: Tuberías de entrenamiento e inferencia inspiradas en SBERT que soportan puntuación de similitud, recuperación, agrupamiento y exportación persistente de incrustaciones.
    \item \textbf{Interfaz de Configuración Interactiva}: Servidor web basado en Streamlit que proporciona generación de formularios impulsada por esquemas, validación en tiempo real, documentación en línea y generación de comandos CLI.
    \item \textbf{Tubería de Validación Automatizada}: Integración continua y pruebas unitarias ampliadas que verifican las rutas de entrenamiento del modelo, la integración optimizador/esquema y la compatibilidad de ejemplos YAML.
\end{enumerate}

\subsection{Guía de Lectura}

Este documento está organizado como una referencia técnica que aborda cuatro preocupaciones operativas:
\begin{enumerate}[leftmargin=*]
    \item \textbf{Requisitos Previos}: El Apéndice~\ref{sec:annex-tutorial} proporciona una introducción accesible a los transformers y mecanismos de atención para lectores nuevos en el campo.
    \item \textbf{Contrato de Configuración}: La Sección~\ref{sec:config-centric-architecture} describe el esquema YAML que garantiza experimentos válidos y la Sección~\ref{sec:schema-validation} explica las reglas de validación.
    \item \textbf{Selección de Arquitectura}: La Sección~\ref{sec:attention-families} proporciona una comparación exhaustiva de las familias de mezcladores de secuencia; los Apéndices~\ref{sec:annex-transformer}, \ref{sec:annex-latent}, \ref{sec:annex-sparse} y \ref{sec:annex-gated} sintetizan la literatura de respaldo; el Apéndice~\ref{sec:annex-norm-bitnet-mod} detalla las variantes de normalización, la cuantización BitNet, el enrutamiento Mixture-of-Depths y la memoria condicional Engram.
    \item \textbf{Dinámica de Optimización}: La Sección~\ref{sec:optimizer-families} detalla el enrutamiento de optimizadores y la dinámica de entrenamiento; el Apéndice~\ref{sec:annex-optimizers} proporciona un análisis exhaustivo de las familias de optimizadores.
    \item \textbf{Despliegue e Inferencia}: Las Secciones~\ref{sec:quantization} y~\ref{sec:sbert} describen el despliegue cuantizado y los flujos de trabajo de incrustación de oraciones.
\end{enumerate}

\subsection{Constructor de Configuración Basado en Web}

Además de la edición directa de YAML, este proyecto proporciona una interfaz web basada en Streamlit (accesible mediante el comando \texttt{web-server}) que mejora la accesibilidad y descubribilidad de la configuración. La interfaz presenta campos de esquema con:

\begin{itemize}[leftmargin=*]
    \item \textbf{Renderizado de formularios impulsado por esquemas} --- Todos los campos se generan dinámicamente a partir del esquema autoritativo, garantizando consistencia y validación.
    \item \textbf{Documentación de parámetros en línea} --- Cada campo del formulario muestra un \emph{título} del esquema como su etiqueta, con la \emph{descripción} mostrada como información sobre herramientas al pasar el cursor.
    \item \textbf{Vista previa de configuración en tiempo real} --- Los usuarios ven la salida YAML en vivo mientras modifican los campos del formulario, permitiendo retroalimentación de validación inmediata.
    \item \textbf{Generación de comandos CLI} --- La interfaz genera comandos CLI completos para entrenamiento, despliegue, inferencia y flujos de trabajo SBERT basados en la configuración actual.
    \item \textbf{Mejoras de accesibilidad} --- La información sobre herramientas y los formularios estructurados facilitan la comprensión de las opciones de configuración, especialmente para usuarios nuevos en el proyecto o que exploran arquitecturas novedosas.
\end{itemize}

Este enfoque basado en web aborda barreras de usabilidad comunes en la experimentación impulsada por configuración:
\begin{itemize}[leftmargin=*]
    \item Reduce la necesidad de memorizar la estructura YAML y los nombres de los campos
    \item Previene errores tipográficos mediante la validación del esquema
    \item Proporciona contexto educativo a través de documentación en línea
    \item Permite la experimentación rápida con ajuste guiado de parámetros
    \item Sirve tanto como herramienta de configuración como recurso de aprendizaje
\end{itemize}

La implementación de la interfaz web utiliza widgets de formulario de Streamlit con:
\begin{itemize}[leftmargin=*]
    \item \texttt{st.checkbox()} para alternativas binarias con texto de ayuda
    \item \texttt{st.number\_input()} para campos numéricos con tamaño de paso y formato
    \item \texttt{st.selectbox()} para opciones de enumeración con visualización de opciones
    \item \texttt{st.multiselect()} para selecciones de arreglo a partir de opciones definidas
    \item \texttt{st.text\_input()} para campos de cadena
    \item \texttt{st.info()}, \texttt{st.caption()} para información complementaria
\end{itemize}

Los metadatos del esquema (campos de título y descripción) se extraen y renderizan sistemáticamente en todas las secciones del formulario, incluyendo:
\begin{itemize}[leftmargin=*]
    \item Parámetros de arquitectura del modelo (tamaño oculto, capas, cabezas de atención, etc.)
    \item Configuraciones de ejecución de entrenamiento (tamaño de lote, acumulación, programador)
    \item Configuración del optimizador con hiperparámetros por grupo de parámetros
    \item Opciones de despliegue y cuantización
    \item Parámetros específicos de entrenamiento e inferencia SBERT
\end{itemize}

